% ============================================================================
% LANGENTWURF LE-008: Passwörter
% Profil: S (Senioren 65+) | Tier: 2 (Standard)
% ============================================================================

\documentclass[11pt,a4paper]{article}
\usepackage[utf8]{inputenc}
\usepackage[T1]{fontenc}
\usepackage[ngerman]{babel}
\usepackage{geometry}
\usepackage{fancyhdr}
\usepackage{lastpage}
\usepackage{xcolor}
\usepackage{tcolorbox}
\usepackage{enumitem}
\usepackage{longtable}
\usepackage{hyperref}
\usepackage{amssymb}

\geometry{left=2.5cm, right=2.5cm, top=2.5cm, bottom=2.5cm}
\definecolor{fach}{HTML}{b35c00}
\definecolor{gray}{HTML}{666666}
\definecolor{successgreen}{HTML}{2a7a4b}
\definecolor{warnred}{HTML}{c62828}

\tcbuselibrary{skins,breakable}
\newtcolorbox{scorebox}{ colback=white, colframe=fach, fonttitle=\bfseries, title=Didaktik-Score, breakable }
\newtcolorbox{warnbox}[1][]{ colback=warnred!5, colframe=warnred, fonttitle=\bfseries, title=#1, breakable }

\pagestyle{fancy}
\fancyhf{}
\fancyhead[L]{\small\textcolor{gray}{Digitale Grundlagen -- 008 Passwörter}}
\fancyhead[R]{\small\textcolor{gray}{LE Tier 2 / Profil S}}
\fancyfoot[C]{\small\thepage{} / \pageref{LastPage}}
\renewcommand{\headrulewidth}{0.4pt}

\begin{document}

\begin{titlepage}
    \centering
    \vspace*{2cm}
    {\Large\textcolor{gray}{Digitale Grundlagen -- Senioren 65+}}\\[2cm]
    {\Huge\textcolor{fach}{\textbf{Langentwurf}}}\\[0.5cm]
    {\large Tier 2 | Profil S}\\[2cm]
    {\LARGE\textbf{008 -- Passwörter verstehen \& verwalten}}\\[0.5cm]
    {\large Modul B: Verbindungen \& Sicherheit}\\[2cm]
    \begin{tabular}{rl}
        \textbf{Zeitrahmen:} & 75--90 Minuten \\
        \textbf{Vorkenntnisse:} & Apple-ID kennen, Einstellungen öffnen \\
    \end{tabular}
    \vfill
    {\normalsize Stand: \today}
\end{titlepage}

\tableofcontents
\newpage

\section{Bedingungsanalyse}

\subsection{Ausgangssituation}
\begin{itemize}
    \item Passwörter sind häufigste Frustquelle bei Senioren
    \item Oft nur ein Passwort für alles (unsicher!)
    \item Passwörter werden auf Zetteln notiert (verloren, gefunden)
    \item iCloud-Schlüsselbund ist unbekannt
    \item Passwort-Wiederherstellung ist Stressfaktor
\end{itemize}

\subsection{Typische Probleme}
\begin{itemize}
    \item ``Ich vergesse ständig meine Passwörter''
    \item ``Ich habe überall das gleiche Passwort''
    \item ``Wo soll ich die alle aufschreiben?''
    \item ``Das Gerät fragt, ob es das Passwort speichern soll -- ist das sicher?''
\end{itemize}

\section{Sachanalyse}

\subsection{Kernkonzepte}
\begin{enumerate}
    \item \textbf{Warum verschiedene Passwörter?} Ein Leck = nur ein Konto betroffen
    \item \textbf{Gutes Passwort:} Lang, Mischung, keine persönlichen Daten
    \item \textbf{iCloud-Schlüsselbund:} Apple speichert Passwörter sicher
    \item \textbf{Automatisches Ausfüllen:} Face ID + gespeichertes Passwort
    \item \textbf{Passwort-Wiederherstellung:} E-Mail, Sicherheitsfragen, Vertrauenskontakt
\end{enumerate}

\subsection{Alltagsanalogien}
\begin{tabular}{|l|p{8cm}|}
\hline
\textbf{Begriff} & \textbf{Analogie} \\
\hline
Passwort & Schlüssel zu einem Schließfach \\
Verschiedene Passwörter & Jedes Schließfach hat eigenen Schlüssel \\
Schlüsselbund & Der Butler, der alle Schlüssel sicher verwahrt \\
Automatisches Ausfüllen & Der Butler öffnet für dich, wenn er dein Gesicht erkennt \\
Passwort-Wiederherstellung & Der Ersatzschlüssel beim Hausmeister \\
\hline
\end{tabular}

\section{Didaktische Analyse}

\subsection{Sicherheitsrelevanz}
\begin{warnbox}[Sicherheitskritisch]
\begin{itemize}
    \item Gleiches Passwort überall = Alle Konten in Gefahr bei einem Leck
    \item Einfache Passwörter (123456, Geburtsdatum) werden in Sekunden geknackt
    \item Passwörter auf Zetteln können gefunden werden
\end{itemize}
\end{warnbox}

\subsection{Didaktische Reduktion}
\textbf{Ausgeklammert:} Passwort-Manager von Drittanbietern, Zwei-Faktor für alle Dienste, Passkeys.

\textbf{Fokus:} iCloud-Schlüsselbund als sichere, einfache Lösung.

\section{Lernziele}

\begin{tabular}{|c|p{7cm}|c|c|}
\hline
\textbf{Nr.} & \textbf{Die Teilnehmer können...} & \textbf{Stufe} & \textbf{Niveau} \\
\hline
FZ1 & erklären, warum verschiedene Passwörter wichtig sind & Verstehen & Basis \\
FZ2 & ein sicheres Passwort nach Merksatz-Methode erstellen & Anwenden & Basis \\
FZ3 & ihre gespeicherten Passwörter im Schlüsselbund finden & Anwenden & Basis \\
FZ4 & das automatische Ausfüllen nutzen & Anwenden & Basis \\
FZ5 & ein vergessenes Passwort zurücksetzen & Anwenden & Vertiefung \\
\hline
\end{tabular}

\section{Methodische Analyse}

\subsection{Merksatz-Methode}
Sichere Passwörter durch Merksätze:
\begin{enumerate}
    \item Satz bilden: ``Meine Enkelin Marie ist 2019 geboren!''
    \item Anfangsbuchstaben: MEMi2019g!
    \item Variieren für verschiedene Dienste: MEMi2019g!-Bank, MEMi2019g!-Email
\end{enumerate}

\section{Verlaufsplanung}

\begin{longtable}{|p{1cm}|p{2cm}|p{5cm}|p{3cm}|p{2cm}|}
\hline
\textbf{Zeit} & \textbf{Phase} & \textbf{Inhalt} & \textbf{TN-Aktivität} & \textbf{Material} \\
\hline
0--10 & Einstieg & ``Wie viele Passwörter haben Sie?'' -- Problem aufzeigen & Nachdenken & -- \\
\hline
10--25 & Theorie & Warum verschiedene? Was ist sicher? Merksatz-Methode & Zuhören, Merksatz erstellen & S.1 \\
\hline
25--40 & Schlüsselbund & Einstellungen $\rightarrow$ Passwörter zeigen, Passwort nachschlagen & Mitmachen & S.2 \\
\hline
40--55 & Automatisch & Automatisches Ausfüllen in Safari demonstrieren & Beobachten, ggf. testen & S.2 \\
\hline
55--70 & Wiederherstellung & Was tun bei vergessenem Passwort? Schritte zeigen & Notieren & S.2 \\
\hline
70--85 & Übungen & Merksatz üben, Checkliste, Notfall-Kontakte & Bearbeiten & S.3 \\
\hline
\end{longtable}

\section{Lösungen}

\textbf{Wo finde ich meine gespeicherten Passwörter?}
\begin{itemize}
    \item iPhone/iPad: Einstellungen $\rightarrow$ Passwörter (Face ID zum Öffnen)
    \item Mac: Systemeinstellungen $\rightarrow$ Passwörter
    \item Safari: Safari $\rightarrow$ Einstellungen $\rightarrow$ Passwörter
\end{itemize}

\textbf{Beispiel-Merksätze:}
\begin{itemize}
    \item ``Jeden Morgen trinke ich 2 Tassen Kaffee!'' $\rightarrow$ JMti2TK!
    \item ``Mein Hund Bello ist 7 Jahre alt.'' $\rightarrow$ MHBi7Ja.
\end{itemize}

\section{Didaktik-Score}

\begin{scorebox}
\begin{tabular}{|c|l|c|c|p{3.5cm}|}
\hline
\# & Dimension & Gew. & Score & Notiz \\
\hline
1 & Differenzierung & 15\% & 8/10 & Basis klar, Vertiefung gut \\
2 & Taxonomie & 10\% & 9/10 & Verstehen + Anwenden \\
3 & Alltagsrelevanz & 10\% & 10/10 & Fundamentales Problem \\
4 & Aktivierung & 10\% & 9/10 & Merksatz selbst erstellen \\
5 & Scaffolding & 10\% & 10/10 & Merksatz-Methode exzellent \\
6 & Feedback & 10\% & 8/10 & Checkliste vorhanden \\
7 & Sprachliche Klarheit & 10\% & 9/10 & Butler-Analogie hilft \\
8 & Motivation & 10\% & 10/10 & Sicherheit motiviert \\
9 & Barrierefreiheit & 10\% & 9/10 & Große Darstellung \\
10 & Praktikabilität & 5\% & 8/10 & 85 Min evtl. lang \\
\hline
\multicolumn{3}{|r|}{\textbf{GESAMT:}} & \textbf{9.1/10} & \textcolor{successgreen}{\textbf{FREIGABE}} \\
\hline
\end{tabular}
\end{scorebox}

\end{document}
